%%% FIELD proof-legal-location-embedding
Fix a sufficiently large even \(n\), condition on \(\mathcal F_n\), and write
\[
 H=H_n,\qquad \boldsymbol u=\boldsymbol 1_n,\qquad
 \boldsymbol v=\boldsymbol v_n,\qquad h_\ell=(-1)^\ell .
\]
The definition of \(\boldsymbol v_n\) gives
\[
 \boldsymbol u^{\top}\boldsymbol v=0,
 \qquad \|\boldsymbol u\|_2^2=\|\boldsymbol v\|_2^2=n.
\]
Consequently
\[
 P:=\frac{\boldsymbol u\boldsymbol u^{\top}
              +\boldsymbol v\boldsymbol v^{\top}}{n}
\]
is the orthogonal projector onto
\(\operatorname{span}\{\boldsymbol u,\boldsymbol v\}\).  By
\(\mathrm{def{:}legal\mbox{-}location\mbox{-}submodel}\), every population
design matrix appearing in that construction is
\[
 B_n=\boldsymbol u\boldsymbol u^{\top}
       +\boldsymbol v\boldsymbol v^{\top}=nP.
\]
Thus \(B_n\) has rank two, both nonzero singular values are \(n\), and
\[
 B_n^\dagger=\frac{P}{n}
 =\frac{\boldsymbol u\boldsymbol u^{\top}
          +\boldsymbol v\boldsymbol v^{\top}}{n^2}.
\]
The definitions of \(\beta_n\), \(a_{n,\ell}\), and \(a_n\) therefore give,
respectively,
\[
 \beta_n=\frac{\cos(\varphi)\boldsymbol u
                     +\sin(\varphi)\boldsymbol v}{n},
 \qquad
 a_{n,\ell}=\frac{\boldsymbol u+h_\ell\boldsymbol v}{n},
 \qquad
 a_n=\frac{\boldsymbol u-\boldsymbol v}{n}.
\]
In particular,
\[
 \|\beta_n\|_2^2=\frac1n,
 \qquad
 a_{n,\ell}^{\top}a_{n,k}=\frac{1+h_\ell h_k}{n},
 \qquad
 \|a_n\|_2^2=\frac2n.                                      \tag{1}
\]

We next verify every member condition in
\(\mathrm{def{:}panel\mbox{-}class}\).  The constant and alternating temporal
factor coordinates, paired with the row loading \((1,(-1)^i)^{\top}\), give
\[
 B_n(i,j)=1+(-1)^{i+j}.
\]
Pairing the same temporal factors with the target loading
\((\cos(\varphi),\sin(\varphi))^{\top}\) gives
\[
 \bar y_n(\varphi)=\cos(\varphi)\boldsymbol u
                    +\sin(\varphi)\boldsymbol v,
 \qquad
 \theta_{n,\ell}=\cos(\varphi)+h_\ell\sin(\varphi).          \tag{2}
\]
This explicitly realizes \(\mathrm{ass{:}factor\mbox{-}model}\).  A Hankel
matrix formed from the constant sequence is an outer product of two all-one
vectors, and a Hankel matrix formed from the alternating sequence is an outer
product of two sign vectors.  Each therefore has rank one, verifying
\(\mathrm{ass{:}hankel\mbox{-}rank}\) with \(Q_0\geq1\).

All five dimensions in the construction equal \(n\), which verifies
\(\mathrm{ass{:}balanced\mbox{-}dimensions}\).  The preceding singular-value
calculation verifies \(\mathrm{ass{:}fixed\mbox{-}rank}\) and
\(\mathrm{ass{:}strong\mbox{-}spectrum}\) with fixed admissible envelopes
containing rank two and singular value \(n\).  Every displayed population
entry, response mean, and target mean has absolute value at most \(2\), so
\(\mathrm{ass{:}bounded\mbox{-}means}\) holds with any declared envelope at
least \(2\).  Equivalently, the fixed radii can be chosen as
\(r_0=2\), \(c=C=c_s=C_s=1\), and \(C_\mu=2\), or loosened without affecting
the construction.

Because
\[
 \operatorname{col}(B_n)=\operatorname{span}\{\boldsymbol u,\boldsymbol v\}
 =\operatorname{col}(B_n^{\top}),                             \tag{3}
\]
\(\mathrm{ass{:}unit\mbox{-}recovery}\) holds.  Furthermore,
\[
 B_na_{n,\ell}=\boldsymbol u+h_\ell\boldsymbol v,
 \qquad B_na_n=\boldsymbol u-\boldsymbol v,                   \tag{4}
\]
so (3)--(4) verify \(\mathrm{ass{:}direct\mbox{-}response\mbox{-}span}\).
The direct transport equality required by
\(\mathrm{ass{:}direct\mbox{-}terminal\mbox{-}transport}\) follows by direct
calculation, rather than by stipulation:
\[
\begin{aligned}
 a_{n,\ell}^{\top}\bar W_n^{\top}\beta_n
 &=a_{n,\ell}^{\top}B_n\beta_n\\
 &=\frac{(\boldsymbol u+h_\ell\boldsymbol v)^{\top}
              (\cos(\varphi)\boldsymbol u+sin(\varphi)\boldsymbol v)}{n}\\
 &=\cos(\varphi)+h_\ell\sin(\varphi)
 =\theta_{n,\ell},
\end{aligned}                                                  \tag{5}
\]
where the last equality is (2).

It remains to check recursive recovery and stability.  Let
\(A=\Pi_n(a_n)=S_n+e_na_n^{\top}\), where \(e_n=e_{L_n}\) because
\(L_n=n\).  Since \(n\) is even,
\[
 a_n^{\top}\boldsymbol u=1,
 \qquad a_n^{\top}\boldsymbol v=-1,
 \qquad S_n\boldsymbol u=\boldsymbol u-e_n,
 \qquad S_n\boldsymbol v=-\boldsymbol v+e_n.
\]
Hence
\[
 A\boldsymbol u=\boldsymbol u,
 \qquad A\boldsymbol v=-\boldsymbol v.                       \tag{6}
\]
Every population terminal state and every one of its iterates under this valid
recurrence consequently remains in the space on the right side of (3), which
is \(\operatorname{col}(\bar Z_{\mathrm{rec},n})\).  This verifies
\(\mathrm{ass{:}recursive\mbox{-}recovery}\).  Moreover,
\[
 \|A\|_{\mathrm{op}}
 \leq \|S_n\|_{\mathrm{op}}+\|e_na_n^{\top}\|_{\mathrm{op}}
 =1+\sqrt{2/n}.
\]
For \(0\leq j\leq H=\lfloor n^{1/8}\rfloor\), submultiplicativity and
\(1+x\leq e^x\) imply
\[
 \|A^j\|_{\mathrm{op}}
 \leq(1+\sqrt{2/n})^j
 \leq\exp(\sqrt2 H/\sqrt n)
 \leq e^{\sqrt2}.                                             \tag{7}
\]
Thus \(\mathrm{ass{:}stable\mbox{-}short\mbox{-}recursion}\) holds uniformly.

The menus are singletons by the construction, so their measurable selections
are deterministic.  The mutually disjoint observed blocks in
\(\mathrm{def{:}legal\mbox{-}location\mbox{-}submodel}\), together with
\(\mathrm{ass{:}gaussian\mbox{-}cells}\), verify
\(\mathrm{ass{:}independent\mbox{-}selection}\) and
\(\mathrm{ass{:}disjoint\mbox{-}inference\mbox{-}blocks}\).  The same Gaussian
condition gives independent \(N(0,\sigma_n^2)\) primitive cell errors, and
\(\mathrm{ass{:}noise\mbox{-}scale}\) supplies
\(0<c_\sigma\leq\sigma_n\leq C_\sigma<\infty\).

We now calculate the direct covariance and verify the remaining variance
floor.  The Riesz definitions, (3), and the fact that the displayed
coefficients belong to \(\operatorname{col}(P)\) yield
\[
 p_{\mathrm{dir},n,\ell}
 =B_n^{\dagger\top}B_na_{n,\ell}=Pa_{n,\ell}=a_{n,\ell},
 \qquad
 q_{\mathrm{dir},n,\ell}
 =B_n^{\dagger\top}B_n^{\top}\beta_n=P\beta_n=\beta_n.       \tag{8}
\]
For later use in the recursive variance bound, put
\(g_{n,\ell}=(A^\ell)^{\top}e_n\).  Equation (6) implies
\[
 \boldsymbol u^{\top}g_{n,\ell}=e_n^{\top}A^\ell\boldsymbol u=1,
 \qquad
 \boldsymbol v^{\top}g_{n,\ell}=e_n^{\top}A^\ell\boldsymbol v=h_\ell.
\]
Therefore
\[
 p_{\mathrm{rec},n,\ell}
 =B_n^{\dagger\top}B_ng_{n,\ell}
 =Pg_{n,\ell}
 =\frac{\boldsymbol u+h_\ell\boldsymbol v}{n},
 \qquad
 \|p_{\mathrm{rec},n,\ell}\|_2^2=\frac2n.                  \tag{9}
\]
In the recursive score of
\(\mathrm{def{:}stacked\mbox{-}scores}\), the term
\(p_{\mathrm{rec},n,\ell}^{\top}e_{y,n}\) is independent of all other
primitive-block terms.  Hence variances add, and (9) gives
\[
 s_{\mathrm{rec},n,\ell}^2
 \geq \sigma_n^2\|p_{\mathrm{rec},n,\ell}\|_2^2
 =\frac{2\sigma_n^2}{n}.                                    \tag{10}
\]

For an i.i.d. mean-zero Gaussian matrix \(E\) with cell variance
\(\sigma_n^2\), direct expansion of the double sum gives
\[
 \operatorname{Cov}(x^{\top}Ey,x'^{\top}Ey'\mid\mathcal F_n)
 =\sigma_n^2(x^{\top}x')(y^{\top}y').                       \tag{11}
\]
Apply (11) separately to the mutually independent terminal, unit-regression,
and temporal-regression blocks in the direct score.  Distinct direct response
errors are independent, while the same \(E_{Z,\mathrm{dir},n}\) is shared
across leads.  Using (8), we obtain
\[
\begin{aligned}
 \frac{\operatorname{Cov}(G_{\mathrm{dir},n,\ell},
                           G_{\mathrm{dir},n,k}\mid\mathcal F_n)}{\sigma_n^2}
 &=\|\beta_n\|_2^2a_{n,\ell}^{\top}a_{n,k}
 +(1+\|\beta_n\|_2^2)
       p_{\mathrm{dir},n,\ell}^{\top}p_{\mathrm{dir},n,k}\\
 &\quad+\|\beta_n\|_2^2
       \{\boldsymbol 1\{\ell=k\}+a_{n,\ell}^{\top}a_{n,k}\}.
\end{aligned}                                                 \tag{12}
\]
Substituting (1) and (8) into (12) gives the complete covariance formula
\[
 \frac{\Sigma_n^{\mathrm{loc}}(\ell,k)}{\sigma_n^2}
 =\left(1+\frac3n\right)\frac{1+h_\ell h_k}{n}
  +\frac{\boldsymbol 1\{\ell=k\}}{n}.                      \tag{13}
\]
It is independent of \(\varphi\), and its diagonal is
\[
 s_{\mathrm{dir},n,\ell}^2
 =\sigma_n^2\left(\frac3n+\frac6{n^2}\right).              \tag{14}
\]
Equations (10), (14), and \(\sigma_n\geq c_\sigma\) verify
\(\mathrm{ass{:}variance\mbox{-}floor}\), for example with any
\(c_v\leq\sqrt2c_\sigma\).

The last summand on the right side of (12) containing
\(\boldsymbol 1\{\ell=k\}\) is exactly the covariance of the direct
lead-response terms
\(q_{\mathrm{dir},n,\ell}^{\top}e_{\mathrm{dir},n,\ell}\).  Their variances
are \(\sigma_n^2\|\beta_n\|_2^2=\sigma_n^2/n\), and the construction makes
them independent across \(\ell\) and independent of all remaining score
terms.  By (14), their standardized variance share is
\[
 \frac{\sigma_n^2/n}
      {\sigma_n^2(3/n+6/n^2)}=\frac1{3+6/n}.                 \tag{15}
\]
In matrix form, (13) is
\[
 \Sigma_n^{\mathrm{loc}}
 =\frac{\sigma_n^2}{n}I_H
  +\frac{\sigma_n^2(1+3/n)}{n}
    (\boldsymbol 1_H\boldsymbol 1_H^{\top}
      +\boldsymbol h_n\boldsymbol h_n^{\top})
 \succeq \frac{\sigma_n^2}{n}I_H.                           \tag{16}
\]
Thus \(\Sigma_n^{\mathrm{loc}}\) is positive definite.  The direct score is
a linear transformation of the primitive Gaussian errors, so it is jointly
Gaussian with mean zero and covariance (13).  From the definition of \(X_n\)
and (2),
\[
 X_n\mid\mathcal F_n
 \sim N(\boldsymbol\vartheta_n(\varphi),\Sigma_n^{\mathrm{loc}}). \tag{17}
\]
For all sufficiently large \(n\), \(H\geq2\), and
\(\boldsymbol 1_H\) and \(\boldsymbol h_n\) are linearly independent.
Consequently the map in (2) is nonconstant (indeed injective on
\([-\pi/6,\pi/6]\)), while (16) makes the Gaussian experiment nondegenerate.
We have checked every defining condition of \(\mathcal C_n\), so
\(\mathcal L_n\subset\mathcal C_n\); equations (14), (15), and (17) prove all
the remaining assertions.

%%% FIELD proof-centered-rectangle-width
Fix \(0<\eta<1\), condition throughout on \(\mathcal F_n\), and abbreviate
\[
 H=H_n,\qquad
 \nu_n=(3+6/n)^{-1},\qquad
 t_n=(1-\eta)\sqrt{2\nu_n\log H}.
\]
By \(\mathrm{prop{:}legal\mbox{-}location\mbox{-}embedding}\), the standardized
direct error
\[
 Z_{n,\ell}
 :=\frac{X_{n,\ell}-\vartheta_\ell}{s_{\mathrm{dir},n,\ell}}
 =\frac{G_{\mathrm{dir},n,\ell}}{s_{\mathrm{dir},n,\ell}}
\]
has an independent direct lead-response component with variance share
\(\nu_n\).  More explicitly, if
\[
 \xi_{n,\ell}
 :=\frac{q_{\mathrm{dir},n,\ell}^{\top}
              e_{\mathrm{dir},n,\ell}}{\sigma_n/\sqrt n},
\]
then the score calculation in that proposition shows that
\(\xi_{n,1},\ldots,\xi_{n,H}\) are independent standard normal variables,
independent of all remaining primitive-block terms.  Hence there is a jointly
Gaussian vector \(\boldsymbol R_n\), independent of
\((\xi_{n,1},\ldots,\xi_{n,H})\), such that
\[
 Z_{n,\ell}=R_{n,\ell}+\sqrt{\nu_n}\,\xi_{n,\ell},
 \qquad 1\leq\ell\leq H.                                    \tag{18}
\]
No independence among the coordinates of \(\boldsymbol R_n\) is required.

We first prove a uniform small-ball bound.  If \(\xi\sim N(0,1)\), then for
every \(x\geq0\) and \(b\in\mathbb R\),
\[
 \Pr\{|\xi-b|<x\}\leq\Pr\{|\xi|<x\}.                       \tag{19}
\]
Indeed, for \(b\geq0\), differentiating
\(\Phi(b+x)-\Phi(b-x)\) with respect to \(b\) gives
\(\phi(b+x)-\phi(b-x)\leq0\); symmetry handles \(b<0\).
Condition on \(\boldsymbol R_n\) in (18), apply (19) coordinatewise, and use
the conditional independence of the \(\xi_{n,\ell}\)'s.  This gives
\[
\begin{aligned}
 &\Pr\left\{\max_{\ell\leq H}|Z_{n,\ell}|<t_n
       \mathrel{\big|}\boldsymbol R_n,\mathcal F_n\right\}\\
 &\qquad\leq
 \left[\Pr\left\{|\xi|<\frac{t_n}{\sqrt{\nu_n}}\right\}\right]^H
 =\{1-2\overline\Phi(x_n)\}^{H},
 \qquad x_n=(1-\eta)\sqrt{2\log H}.                         \tag{20}
\end{aligned}
\]
Because \(H=\lfloor n^{1/8}\rfloor\to\infty\), eventually \(x_n\geq1\).
For every \(x\geq1\), monotonicity of the standard normal density gives the
elementary lower bound
\[
\begin{aligned}
 \overline\Phi(x)
 &=\int_x^\infty\phi(u)\,du
 \geq\int_x^{x+1/x}\phi(u)\,du\\
 &\geq\frac1x\phi(x+1/x)
 \geq\frac{e^{-3/2}}{x}\phi(x).                             \tag{21}
\end{aligned}
\]
Substitution of \(x_n\) into (21) shows that, for a universal constant
\(c>0\),
\[
 2H\overline\Phi(x_n)
 \geq c\frac{H^{1-(1-\eta)^2}}{\sqrt{\log H}}
 =c\frac{H^{2\eta-\eta^2}}{\sqrt{\log H}}
 \longrightarrow\infty.                                    \tag{22}
\]
Since \((1-u)^H\leq e^{-Hu}\) for \(0\leq u\leq1\), integrating (20) over
\(\boldsymbol R_n\) and using (22) yields
\[
 r_n:=\Pr\left\{\max_{\ell\leq H}|Z_{n,\ell}|<t_n
                 \mathrel{\big|}\mathcal F_n\right\}
 \leq\exp\{-2H\overline\Phi(x_n)\}=o(1).                  \tag{23}
\]
The deterministic upper bound in (23) is independent of
\(\boldsymbol\vartheta\), because the conditional error law is the same at
every point of \(\Theta_n^{\mathrm{loc}}\).

Let \(C_n^{\mathrm{rect}}\in\mathfrak R_n\), and fix
\(\boldsymbol\vartheta\in\Theta_n^{\mathrm{loc}}\).  Write
\[
 A_{n,\boldsymbol\vartheta}
 :=\{\boldsymbol\vartheta\in C_n^{\mathrm{rect}}\}.
\]
By \(\mathrm{def{:}centered\mbox{-}rectangles}\), on this coverage event one
has, for every coordinate,
\[
 -W^+_{n,\ell}
 \leq \frac{X_{n,\ell}-\vartheta_\ell}{s_{\mathrm{dir},n,\ell}}
 \leq W^-_{n,\ell}.                                        \tag{24}
\]
Consequently the following inclusion holds pointwise in both the data and the
independent auxiliary randomizer:
\[
 A_{n,\boldsymbol\vartheta}\cap\{W_n^{\max}<t_n\}
 \subseteq\left\{\max_{\ell\leq H}|Z_{n,\ell}|<t_n\right\}. \tag{25}
\]
Taking conditional probabilities in (25) and rearranging gives
\[
\begin{aligned}
 \Pr_{\mathbb P_{\boldsymbol\vartheta,n}}
   \{W_n^{\max}\geq t_n\mid\mathcal F_n\}
 &\geq
 \Pr_{\mathbb P_{\boldsymbol\vartheta,n}}
   (A_{n,\boldsymbol\vartheta}\mid\mathcal F_n)-r_n\\
 &\geq1-\alpha-\delta_n-r_n.                               \tag{26}
\end{aligned}
\]
The last line uses the assumed uniform conditional coverage.  Since the same
\(r_n=o(1)\) works for every \(\boldsymbol\vartheta\), taking the infimum in
(26) proves
\[
 \inf_{\boldsymbol\vartheta\in\Theta_n^{\mathrm{loc}}}
 \Pr_{\mathbb P_{\boldsymbol\vartheta,n}}
 \left\{W_n^{\max}\geq
 (1-\eta)\sqrt{2(3+6/n)^{-1}\log H_n}
 \mathrel{\big|}\mathcal F_n\right\}
 \geq1-\alpha-\delta_n-o(1).                               \tag{27}
\]

Finally, \(W_n^{\max}\geq0\).  The elementary tail-integral inequality
\(\mathbb E[V]\geq t\Pr\{V\geq t\}\), valid for every nonnegative random
variable \(V\) and \(t\geq0\), applied conditionally and combined with (26)
gives
\[
 \inf_{\boldsymbol\vartheta\in\Theta_n^{\mathrm{loc}}}
 \mathbb E_{\mathbb P_{\boldsymbol\vartheta,n}}
 [W_n^{\max}\mid\mathcal F_n]
 \geq(1-\eta)\sqrt{2(3+6/n)^{-1}\log H_n}
       \{1-\alpha-\delta_n-o(1)\}.                          \tag{28}
\]
The inclusion (25) did not require symmetric, deterministic, or nonrandom
widths, so (27)--(28) cover the entire declared class \(\mathfrak R_n\).  It
did use the center \(X_n\) in (24); therefore the proof makes no assertion
about alternative centers or unrestricted estimators, exactly as claimed.
